\documentclass{homework}
\course{Math 4182H}
\author{Jim Fowler}
\hwtitle{Honors Analysis II}
\input{preamble}

\newcommand{\Mtwo}{\mathrm{Mat}_{2 \times 2}}
\newcommand{\GL}{\mathrm{GL}}

\begin{document}
\maketitle

\begin{inspiration}
Granny's \textbf{implicit} belief that everything should get out of her way extended to other witches, very tall trees, and, on occasion, mountains. \\
\byline{Terry Pratchett, \textit{Wyrd Sisters} (not about the curve known as the Witch of Agnesi).}
\end{inspiration}

%%%%%%%%%%%%%%%%%%%%%%%%%%%%%%%%%%%%%%%%%%%%%%%%%%%%%%%%%%%%%%%%
\section{Terminology}

\begin{problem}
State the \textbf{Inverse Function Theorem}.
\end{problem}

\begin{problem}
State the \textbf{Implicit Function Theorem}.
\end{problem}

\begin{problem}
  What is a \textbf{critical point} of a function $f : \R^n \to \R$?
\end{problem}

\begin{problem}\label{regular-value-n-to-one}%
  What is a \textbf{regular value} of a function $f : \R^n \to \R$?
\end{problem}

%%%%%%%%%%%%%%%%%%%%%%%%%%%%%%%%%%%%%%%%%%%%%%%%%%%%%%%%%%%%%%%%
\section{Numericals}

\begin{problem}
  Recall \ref{derivative-of-determinant}.

  Denote the set of $2$-by-$2$ matrices with real entries by $\Mtwo(\R)$.

  Determine the regular values of $\det : \Mtwo(\R) \to \R$.
\end{problem}
  
\begin{problem}
  Recall \ref{derivative-of-inverse}.
  
Let $\GL_2(\R)\subset \Mtwo(\R)$ be the set of invertible $2$-by-$2$ matrices. Consider the map
$\mathrm{inv}:\GL_2(\R)\to \GL_2(\R)$ defined by $\mathrm{inv}(A)=A^{-1}$.

Fix $A_0\in \GL_2(\R)$.  Define
\[
F:\Mtwo(\R)\times \Mtwo(\R)\to \Mtwo(\R),\qquad F(A,B)=AB-I.
\]
Note that $F(A,B)=0$ holds exactly when $B=A^{-1}$.

Use ``implicit differentiation'' in order to compute
the derivative $D(\mathrm{inv})_{A_0}(H)$, i.e., the derivative of the inverse map at $A_0 \in \GL_2(\R)$ and in the direction $H\in \Mtwo(\R)$.
\end{problem}

%%%%%%%%%%%%%%%%%%%%%%%%%%%%%%%%%%%%%%%%%%%%%%%%%%%%%%%%%%%%%%%%
\section{Exploration}

\begin{problem}
Define $\Phi:(0,\infty)\times\R\to\R^2$ by
\[
\Phi(r,\theta)=(r\cos\theta,r\sin\theta).
\]
Compute $\det D\Phi(r,\theta)$.

Your friend says ``so every point of $\R^2\setminus\{0\}$ has a neighborhood on which $(x,y)$ can be written as $(r,\theta)$ uniquely.''

Is this precisely right?
\end{problem}

\pagebreak

\begin{problem}
Consider the elliptic curve
\[
E=\{(x,y)\in\mathbb R^2 : y^2 = x^3 - x\}.
\]
Let $F(x,y)=y^2-(x^3-x)$, so that $E=F^{-1}(0)$.
 
Pick $x_0$ so that $p=(x_0,0)\in E$, i.e., pick $x_0\in\{-1,0,1\}$.

Explain why you \emph{cannot} solve locally for $y$ as a function of $x$ around such a point $p$, but you \emph{can} solve locally for $x$ as a function of $y$.

More precisely, show there is an $\epsilon > 0$ and a $C^1$ function $f:(-\epsilon,\epsilon)\to\mathbb R$ such that $f(0)=x_0$ and, when $-\epsilon < y < \epsilon$, we have $(f(y),y) \in E$. In this case, compute $f'(0)$.
\end{problem}

\begin{problem}
The \textbf{folium of Descartes} is the curve in $\R^2$ carved out by
\[
x^3+y^3=3xy.
\]
Assume $y\neq 0$ and set $t=x/y$.

Express the points on the curve with $y\neq 0$ in parametric form
$(x(t),y(t))$ where $t = x/y$. Compute $x'(t)$. For which values of $t$ is $x'(t)\neq 0$?

Use the Inverse Function Theorem to identify points of the folium where the curve is the graph of a differentiable function $y = y(x)$. Locate the two places where this fails, and explain what happens there.
\end{problem}

\begin{problem}
Fix an integer $k\ge 2$, and consider the map
$P: \Mtwo(\R)\to \Mtwo(\R)$
given by 
$P(X)=X^k$. This is the $k$-th power map.

Show that $P$ is differentiable and compute the derivative of $P$ at $I$. Show that the derivative at $I$ is an isomorphism. Conclude that for every sufficiently small matrix $H$, the equation $X^k=I+H$ has a solution.

Is it unique? Does every matrix have a $k$-th root?
\end{problem}

\begin{problem}
  Interpret $\mathbb{R}^{n+1}$ as the coefficients of a degree $n$ polynomial. Pick $P \in \mathbb{R}^{n+1}$ so that there is a root $r$ of the polynomial $P$, i.e.,
  \[
    P_0 + P_1 \, r + P_2 \, r^2 + \cdots + P_n r^n = 0.
  \]
  Does there exist a function $f : U \to \R$ defined on a neighborhood $U \ni P$ so that $f(Q)$ is a root of the polynomial $Q$? In this case, compute $f'(P)$.
\end{problem}

%%%%%%%%%%%%%%%%%%%%%%%%%%%%%%%%%%%%%%%%%%%%%%%%%%%%%%%%%%%%%%%%
\section{Prove or Disprove and Salvage if Possible}

\begin{problem}
Let $F:\R\to\R$ be $C^1$. If $F'(x) \neq 0$ for all $x\in\R$, then $F$ is surjective.
\end{problem}

\begin{problem}
Let $F:\R^n\to\R^n$ be $C^1$. If $\det DF(x)\neq 0$ for all $x\in\R^n$, then $F$ is injective.
\end{problem}

\begin{problem}
  Recall \ref{cauchy-riemann-equations}.

  Suppose $F:\R^2 \to\R^2$ is $C^1$ and $F(x,y) = (u(x,y),v(x,y))$ for functions $u, v : \R^2 \to \R$.  Suppose $F$ satisfies the Cauchy--Riemann equations, i.e., $u_x=v_y$ and $u_y=-v_x$. Assume $DF(x_0,y_0)$ is invertible at some $(x_0,y_0)$, and let $G=(p,q)$ be the $C^1$ local inverse of $F$ near $(x_0,y_0)$. Then $G$ also satisfies the Cauchy--Riemann equations.
\end{problem}

\end{document}

